\begin{center}
\Huge
Differentialligninger og eksponentiel vækst
\end{center}


\section*{Løsninger til differentialligninger}
\stepcounter{section}

Som tidligere nævnt, så bruges differentialligninger ofte til at modellere virkelige fænomener. Vi skal derfor se på forskellige situationer, hvor vi kan gøre visse antagelser om væksten af en model og ud fra dette løse en differentialligning for at beskrive modellen eksplicit. Vi begynder med tilfældet, at væksten af modellen er konstant eller mere præcist, at der for modellen $y$ gælder, at
\begin{align*}
	y' = k.
\end{align*}
Vi har tidligere set, hvordan løsningerne for sådanne differentialligninger ser ud, men vi vil nu være mere præcise.
\begin{setn}
	Differentialligningen 
	\begin{align*}
		y' = k
	\end{align*}
	har den fuldstændige løsning
	\begin{align*}
		y(x) = kx + c,
	\end{align*}
	hvor $c \in \mathbb{R}$. 
\end{setn}
\begin{proof}
	Vi integrerer begge sider af lighedstegnet og får, at 
	\begin{align*}
		y(x) = \int y' d x = \int k d x = kx + c.
	\end{align*}
\end{proof}

En mere interessant model er en model, hvor væksten af modellen $y$ er proportional med $y$ eller mere præcist
\begin{align*}
	y' = ky.
\end{align*}
Vi har tidligere arbejdet med sådanne differentialligninger, men vi skal nu bevise løsningsformlen for denne type differentialligninger. 
Vi kan tegne et hældningsfelt for en sådan type differentialligning. Dette kan ses på Figur \ref{fig:eksphæld}.
\begin{figure}[H]
	\centering
	\begin{tikzpicture}
	\begin{axis}[
		xmin = -5, xmax = 5,
		ymin = -0.5, ymax = 5,
		axis lines = center,
		xlabel = $x$,
		ylabel = $y$,
		ticks = none
	]
		
		\foreach \xVal in {-5,-4.5,...,5}{
			\foreach \yVal in {0,0.25,...,5}{
				\pgfmathsetmacro{\sizescale}{0.1};
				\pgfmathsetmacro{\slope}{ 0.5*\yVal};
				\pgfmathsetmacro{\angle}{atan(\slope)};
				\pgfmathsetmacro{\xlen}{\sizescale*cos(\angle)};
				\pgfmathsetmacro{\ylen}{\sizescale*sin(\angle)};
				\edef\temp{\noexpand\draw [color = SeaGreen] 
				(axis cs: \xVal-\xlen, \yVal-\ylen) -- (axis cs: \xVal + \xlen, \yVal + \ylen);}
				\temp
			}
		}
	\end{axis}
	\end{tikzpicture}
	\caption{Hældningsfelt for differentialligningen $y' = ky$.}
	\label{fig:eksphæld}
\end{figure}
Løsningskurverne kunne godt ud fra hældningsfeltet se ud til at være eksponentielle. Tilmed ved vi jo allerede, at der for den naturlige eksponentialfunktion gælder, at differentialkvotienten og funktionsværdien tilsvarer hinanden. Vores formodning skal da også vise sig at være korrekt ifølge følgende sætning. 
\begin{setn}\label{setn:2}
	Differentialligningen 
	\begin{align*}
		y' = ky
	\end{align*}
	har den fuldstændige løsning
	\begin{align*}
		y(x) = ce^{kx},
	\end{align*}
	hvor $c \in \mathbb{R}$. 
\end{setn}
\begin{proof}
	Beviset består af to dele - Vi skal vise, at $y(x) = ce^{kx}$ er en løsning, og så skal vi vise, at alle løsninger er på formen $y(x) = ce^{kx}$. Første del er en opgave,
	og anden del klarer vi her. 
	Antag derfor, at $y(x)$ er en løsning til $y' = ky$. Vi definerer så en hjælpefunktion
	\begin{align*}
		z(x) = y(x)e^{-kx},
	\end{align*}
	som vi differentierer.
	\begin{align*}
		z'(x) = y'(x)e^{-kx}+y(x)(-k)e^{-kx} = ky(x)e^{-kx}-ky(x)e^{-kx}=0.
	\end{align*}
	Da $z'(x)=0$, så betyder det, at $z(x) = c$. Derfor ved vi, at
	\begin{align*}
		z(x) = y(x)e^{-kx} = c
	\end{align*}
	eller 
	\begin{align*}
		y(x) = ce^{kx},
	\end{align*}
	hvilket var hvad, vi skulle vise. 		 
\end{proof}
Vi kan også tegne en løsningskurve for denne differentialligning ind i hældningsfeltet. Dette kan ses på Figur \ref{fig:eksphældkurve}.
\begin{figure}[H]
	\centering
	\begin{tikzpicture}
	\begin{axis}[
		xmin = -5, xmax = 5,
		ymin = -0.5, ymax = 5,
		axis lines = center,
		xlabel = $x$,
		ylabel = $y$,
		ticks = none
	]
		\addplot[color = Melon, thick, samples = 200, domain = -5:5] {2*exp(0.5*x)};
		\foreach \xVal in {-5,-4.5,...,5}{
			\foreach \yVal in {0,0.25,...,5}{
				\pgfmathsetmacro{\sizescale}{0.1};
				\pgfmathsetmacro{\slope}{ 0.5*\yVal};
				\pgfmathsetmacro{\angle}{atan(\slope)};
				\pgfmathsetmacro{\xlen}{\sizescale*cos(\angle)};
				\pgfmathsetmacro{\ylen}{\sizescale*sin(\angle)};
				\edef\temp{\noexpand\draw [color = SeaGreen] 
				(axis cs: \xVal-\xlen, \yVal-\ylen) -- (axis cs: \xVal + \xlen, \yVal + \ylen);}
				\temp
			}
		}
	\end{axis}
	\end{tikzpicture}
	\caption{Hældningsfelt for differentialligningen $y' = ky$ samt løsningskurve.}
	\label{fig:eksphældkurve}
\end{figure}
\begin{exa}
	En differentialligning er givet ved 
	\begin{align*}
		y' = 10y.
	\end{align*}
	Så er den fuldstændige løsning til differentialligningen givet ved
	\begin{align*}
		y(x) = ce^{10x}.
	\end{align*}
\end{exa}

Den sidste type differentialligning vi skal se på i dag er den forskudte eksponentielle vækst. Vi antager derfor, at væksten af modellen er lineært aftagende som funktion af modellen - Eller mere præcist
\begin{align*}
y' = b-ay.
\end{align*}
Vi tegner som før et hældningsfelt for en god ordens skyld. Denne kan ses på Figur \ref{fig:ekspforskudthæld}, hvor vi kan se, at løsningskurverne er eksponentielle, men har en asymptote langs $y = b/a$ i stedet for $y = 0$.  
\begin{figure}[H]
	\centering
	\begin{tikzpicture}
	\begin{axis}[
		xmin = -5, xmax = 5,
		ymin = -5, ymax = 5,
		axis lines = center,
		xlabel = $x$,
		ylabel = $y$,
		ticks = none
	]
		
		\foreach \xVal in {-5,-4.5,...,5}{
			\foreach \yVal in {-5,-4.5,...,5}{
				\pgfmathsetmacro{\sizescale}{0.15};
				\pgfmathsetmacro{\slope}{2 - 0.75*\yVal};
				\pgfmathsetmacro{\angle}{atan(\slope)};
				\pgfmathsetmacro{\xlen}{\sizescale*cos(\angle)};
				\pgfmathsetmacro{\ylen}{\sizescale*sin(\angle)};
				\edef\temp{\noexpand\draw [color = SeaGreen] 
				(axis cs: \xVal-\xlen, \yVal-\ylen) -- (axis cs: \xVal + \xlen, \yVal + \ylen);}
				\temp
			}
		}
	\end{axis}
	\end{tikzpicture}
	\caption{Hældningsfelt for differentialligningen $y' = b- ay$.}
	\label{fig:ekspforskudthæld}
\end{figure}

\begin{setn}\label{setn:3}
	Differentialligningen
	\begin{align}\label{eq:diff3}
		y' = b - ay
	\end{align}
	har den fuldstændige løsning
	\begin{align*}
		y(x) = \frac{b}{a} + ce^{-ax},
	\end{align*}
	hvor $c \in \mathbb{R}$.
\end{setn}
\begin{proof}
	Beviset består som før af to dele - Vi skal vise, at $y(x) = \frac{b}{a} +ce^{-ax}$ er en løsning samt at enhver løsning vil være på denne form. Igen skal I selv stå for
	første del af beviset - Vi beviser her anden del.
	
	Antag derfor, at $y(x)$ er en løsning til \eqref{eq:diff3}. Vi introducerer så en hjælpefunktion
	\begin{align*}
		z(x) = b-ay(x). 
	\end{align*}
	Igen differentierer vi hjælpefunktionen. 
	\begin{align*}
		z'(x) = -ay'(x) = -a(b-ay(x)) =-az(x). 
	\end{align*}
	Da $z' = -az$, så har $z$ ifølge Sætning \ref{setn:2} løsningen
	\begin{align*}
		z(x) = \tilde{c}e^{-ax}. 
	\end{align*}
	Derfor fås nu
	\begin{align*}
		b-ay(x) = \tilde{c}e^{-ax} \ &\Leftrightarrow \ -ay(x) = -b+\tilde{c}e^{-ax}\\
		&\Leftrightarrow	\ y(x) = \frac{-b}{-a} + \frac{\tilde{c}}{-a}e^{-ax}\\
		&\Leftrightarrow \ y(x) = \frac{b}{a} + ce^{-ax},
	\end{align*}
	hvor $-\tilde{c}/a = c$. 
\end{proof}
Vi tegner også en løsningskurve i hældningsfeltet. Dette kan ses på Figur \ref{fig:ekspforskudthældkurve}
\begin{figure}[H]
	\centering
	\begin{tikzpicture}
	\begin{axis}[
		xmin = -5, xmax = 5,
		ymin = -5, ymax = 5,
		axis lines = center,
		xlabel = $x$,
		ylabel = $y$,
		ticks = none
	]
		\addplot[color = Melon, thick, domain = -5:5, 
		samples = 200] {2/0.75 - exp(-0.75*x)};
		\foreach \xVal in {-5,-4.5,...,5}{
			\foreach \yVal in {-5,-4.5,...,5}{
				\pgfmathsetmacro{\sizescale}{0.15};
				\pgfmathsetmacro{\slope}{2 - 0.75*\yVal};
				\pgfmathsetmacro{\angle}{atan(\slope)};
				\pgfmathsetmacro{\xlen}{\sizescale*cos(\angle)};
				\pgfmathsetmacro{\ylen}{\sizescale*sin(\angle)};
				\edef\temp{\noexpand\draw [color = SeaGreen] 
				(axis cs: \xVal-\xlen, \yVal-\ylen) -- (axis cs: \xVal + \xlen, \yVal + \ylen);}
				\temp
			}
		}
	\end{axis}
	\end{tikzpicture}
	\caption{Hældningsfelt for differentialligningen $y' = b- ay$ med en løsningskurve.}
	\label{fig:ekspforskudthældkurve}
\end{figure}

\begin{exa}
	Vi skal bestemme en partikulær løsning til differentialligningen	
	\begin{align*}
		y' = 8-5y,
	\end{align*}
	der går gennem punktet $(0,3)$. Ifølge Sætning \ref{setn:3} har differentialligningen den generelle løsning
	\begin{align*}
		y(x) = \frac{8}{5} + ce^{-5x}.
	\end{align*}
	Vi indsætter nu det kendte punkt:
	\begin{align*}
		3=y(0) = \frac{8}{5} + ce^{-5\cdot 0} = \frac{8}{5}+ c. 
	\end{align*}
	Vi får derfor, at $c = 7/5$, og den partikulære løsning lyder da
	\begin{align*}
		y(x) = \frac{8}{5} + \frac{7}{5}e^{-5x}.
	\end{align*}
\end{exa}

\subsection*{Opgave 1}
\begin{enumerate}[label=\roman*)]
	\item Bestem en fuldstændig løsning til differentialligningen
	\begin{align*}
		y' = 5y.
	\end{align*} 
	\item Bestem en fuldstændig løsning til differentialligningen
	\begin{align*}
		y' = -11y.
	\end{align*}
	\item Bestem en partikulær løsning til differentialligningen
	\begin{align*}
		\frac{d y}{d x} = 6y,
	\end{align*}
	der går gennem punktet $(0,4)$
	\item Bestem en partikulær løsning til differentialligningen
	\begin{align*}
		y' = \frac{1}{2}y,
	\end{align*}
	der går gennem punktet $(0,-3)$
	\item Bestem en partikulær løsning til differentialligningen
	\begin{align*}
		y' = y,
	\end{align*}
	der går gennem punktet $(\ln(4),1)$
	\item Bestem en partikulær løsning til differentialligningen
	\begin{align*}
		y' - 3y = 0,
	\end{align*}
	der går gennem punktet $(2,e^8)$
\end{enumerate}
\subsection*{Opgave 2}
\begin{enumerate}[label=\roman*)]
	\item Bestem en fuldstændig løsning til differentialligningen 
	\begin{align*}
		y' = 7-13y
	\end{align*}
	\item Bestem en fuldstændig løsning til differentialligningen 
	\begin{align*}
		y' = 125 + 5y
	\end{align*}
	\item Bestem en fuldstændig løsning til differentialligningen 
	\begin{align*}
		y' -\ln(6)+\sqrt{2}y=0
	\end{align*}
	\item Bestem en partikulær løsning til differentialligningen 
	\begin{align*}
		y' = 8-2y,
	\end{align*}
	der går gennem punktet $(0,10)$. 
	\item Bestem en partikulær løsning til differentialligningen 
	\begin{align*}
		y' = -20 + 5y,
	\end{align*}
	der går gennem punktet $(2,4 + 11e^{10})$. 
\end{enumerate}
\subsection*{Opgave 3}
Vi skal færdiggøre beviserne for Sætningerne \ref{setn:2} og \ref{setn:3}.
\begin{enumerate}[label=\roman*)]
	\item Vis, at 
	\begin{align*}
		y(x) = ce^{kx}
	\end{align*}
	er en løsning til differentialligningen
	\begin{align*}
		y' = ky
	\end{align*}
	\item Vis, at 
	\begin{align*}
		y(x) = \frac{b}{a} + ce^{-ax}
	\end{align*}
	er en løsning til differentialigningen 
	\begin{align*}
		y' = b-ay.
	\end{align*}
\end{enumerate}


\subsection*{Opgave 4 (Med hjælpemidler)}
	Et objekt med en temperatur på $100^\circ$C stilles i et lokale med en temperatur på $19^\circ$C. Det er oplagt, at temperaturen falder langsomme og langsommere jo nærmere
	temperaturen kommer på den omkringliggende temperatur. Vi antager, at temperaturfaldets hastighed er proportional med objektets temperatur eller mere præcist, at 
	temperaturen og temperaturfaldet opfylder differentialligningen
	\begin{align*}
		y' = -k(y-19) = -ky+19k.
	\end{align*}
	\begin{enumerate}[label=\roman*)]
		\item Brug Sætning \ref{setn:3} til at bestemme den fuldstændige løsning til differentialligningen. 
		\item Udnyt, at $y(0) = 100$ til at bestemme konstanten $c$. 
		\item Udnyt, at temperaturen efter 5 minutter er $80^\circ$C til at bestemme $k$. 
		\item Opskriv modellen $y(t)$ for temperaturen af objektet. 
	\end{enumerate}

\subsection*{Opgave 5 (Med hjælpemidler)}
I en bakteriekoloni er der til tidspunktet $t = 0$ 200 mio bakterier. Det antages, at antallet af bakterier $B$ i en afgrænset periode opfylder differentialligningen
\begin{align*}
	\frac{dB}{dt} = k(B - 50),
\end{align*}
hvor $t$ er antallet af forløbne timer.
\begin{enumerate}[label = \roman*)]
	\item Opskriv den fuldstændige løsning for differentialligningen.
	\item Udnyt, at $B(0) = 200$ for at finde $c$.
	\item Udnyt, at antallet af bakterier fordobles på et døgn til at bestemme $k$.
	\item Bestem antallet af bakterier efter 40 timer. 
\end{enumerate}
\subsection*{Opgave 6 (Med hjælpemidler)}
I en by har der i en årrække været fraflytning. Det antages, at antallet af beboere i byen $N$ kan beskrives ved differentialligningen
\begin{align*}
	\frac{dN}{dt} = -k(N - 20),
\end{align*}
hvor $t$ er antallet af år siden begyndelsesåret for målingerne og $N$ er antallet af beboere i byen.
\begin{enumerate}[label=\roman*)]
	\item Opskriv den fuldstændige løsning for differentialligningen. 
\end{enumerate}
Det oplyses, at der til $t = 0$ var $50$ tusinde beboere i byen.
\begin{enumerate}[label=\roman*)]
	\setcounter{enumi}{1}
	\item Brug dette til at bestemme den partikulære løsning til differentialligningen. 
	\item Bestem antallet af beboere i byen efter 20 år hvis $k = 0.02$
	\item Tegn grafen for løsningen og overvej, hvad der sker, når $t \to \infty$.
\end{enumerate}
